\documentclass[12pt,a4paper]{article}
        \makeatletter
        \def\input@path{{../}}
        \makeatother
        \usepackage{almeria-fichas}

        \FichaMeta{Sesion 04}{Coordenadas y localizacion}{Bloque 1 · Mapa y coordenadas}{Leer, escribir y localizar lugares exactos en el mapa usando pares de coordenadas matematicas.}{Procedimientos}{Comprender, Aplicar, Evaluar}

        \begin{document}

        \FichaCabecera

        \begin{materialesbox}
        \begin{itemize}
    \item Ficha impresa
    \item Lapiz
    \item Regla
    \item Plano cuadriculado del aula o de Almeria
\end{itemize}
        \end{materialesbox}

        \begin{pasosbox}
        \begin{enumerate}
    \item Empieza siempre en el origen.
    \item Cuenta primero los pasos horizontales y luego los verticales.
    \item Escribe las coordenadas con parentesis y coma.
    \item Comprueba una segunda vez antes de pasar a la siguiente actividad.
\end{enumerate}
        \end{pasosbox}

        \begin{recordatoriobox}
        \begin{itemize}
    \item La forma correcta es $(x,y)$.
    \item Cambiar el orden cambia el lugar.
    \item Dos lugares distintos no pueden tener el mismo par de coordenadas si el mapa esta bien hecho.
\end{itemize}
        \end{recordatoriobox}

        \begin{apoyobox}
        \begin{itemize}
    \item Rodea la x y la y en cada consigna.
    \item Puedes trazar lineas suaves hasta el punto para no perderte.
    \item Si dudas entre dos valores, cuenta otra vez desde el origen.
\end{itemize}
        \end{apoyobox}

        \BloomSection{comprender}

\BloomExercise{comprender}{1}{Explica el orden}{Explica con tus palabras por que $(4,2)$ y $(2,4)$ no representan el mismo lugar.}{6}

\BloomExercise{comprender}{2}{Lugar exacto}{Escribe una frase para explicar a un compañero como localizar $M(3,5)$ en un plano.}{6}

\BloomSection{aplicar}

\BloomExercise{aplicar}{1}{Busca y escribe}{En un plano inventado, coloca un museo en $(2,4)$, un parque en $(5,1)$, una plaza en $(3,3)$ y una estacion en $(1,2)$.}{7}

\BloomExercise{aplicar}{2}{Ruta de puntos}{Escribe las coordenadas de \textbf{cinco lugares} de una ruta corta por el mapa. Despues numeralos del 1 al 5.}{7}

\BloomSection{evaluar}

\BloomExercise{evaluar}{1}{Revision de localizaciones}{Un equipo ha puesto: playa $(4,2)$, puerto $(4,2)$, museo $(6,1)$. \textbf{Evalua} si su mapa es claro y explica que deberian revisar.}{7}

\BloomExercise{evaluar}{2}{Eleccion del mejor mapa}{Dos mapas muestran los mismos lugares. Uno tiene coordenadas bien escritas y otro mezcla $(x,y)$ con $(y,x)$. \textbf{Evalua} cual es mas fiable y justifica.}{7}

        \begin{evidenciabox}
        \begin{itemize}
    \item Escribir y leer coordenadas sin invertir el orden.
    \item Relacionar puntos del mapa con lugares concretos.
    \item Justificar si una localizacion es correcta o no.
\end{itemize}
        \end{evidenciabox}

        \end{document}
